% !TEX root = ../main.tex
\subsection*{The Standard Machinery}%
\manualtoc{subsection}{\texorpdfstring{The Standard Machinery}{The Standard Machinery}}{studiabile-nav-3001}%

Many properties of Lebesgue integrals can be established through the following four steps:
\begin{enumerate}
\item Prove the property for indicators $\mathbf1_A$, where $A\in\Sigma$.
\item Use linearity to extend it to nonnegative simple functions.
\item Approximate nonnegative measurable functions by increasing simple functions and apply the monotone convergence theorem (MON, Theorem~5.2).
\item Use $f=f^+-f^-$ to extend the result to signed measurable functions whenever the relevant integrals are well-defined.
\end{enumerate}
This is the standard machinery used throughout the course.

\subsection*{A Measure Defined by a Nonnegative Density}%
\manualtoc{subsection}{\texorpdfstring{A Measure Defined by a Nonnegative Density}{A Measure Defined by a Nonnegative Density}}{studiabile-nav-3002}%

Different measures on the same measurable space yield different Lebesgue integrals. Fix $f\in(m\Sigma)^+$ and define, for every $A\in\Sigma$,
\[
\lambda(A):=\int_S f\mathbf1_A\,d\mu=\int_A f\,d\mu.
\]
The following homework proves that $\lambda:\Sigma\to[0,\infty]$ is a measure. In addition, $\lambda$ is absolutely continuous with respect to $\mu$, written $\lambda\ll\mu$, meaning
\[
\mu(N)=0\quad\Longrightarrow\quad\lambda(N)=0,
\qquad\forall N\in\Sigma.
\]

\subsection*{HW 53 --- Countable Additivity of $\lambda$}%
\manualtoc{subsection}{\texorpdfstring{HW 53 --- Countable Additivity of $\lambda$}{HW 53 --- Countable Additivity of lambda}}{studiabile-nav-3003}%

\studyheading{Statement.} Let $f\in(m\Sigma)^+$ and define $\lambda(A)=\int_S f\mathbf1_A\,d\mu$ for every $A\in\Sigma$. Prove that $\lambda$ is countably additive and hence is a measure on $(S,\Sigma)$.

\par\noindent\studyheading{Proof.}\par\nobreak

First, $\lambda(A)\ge0$ for every $A\in\Sigma$, since $f\mathbf1_A\ge0$, and
\[
\lambda(\emptyset)=\int_S f\mathbf1_{\emptyset}\,d\mu
=\int_S0\,d\mu=0.
\]
The equality follows from $\mathbf1_{\emptyset}=0$ and the definition of the positive integral.

Let $(A_n)_{n\in\mathbb N}\subseteq\Sigma$ be pairwise disjoint. Set
\[
U_r:=\bigcup_{n=1}^{r}A_n,
\qquad U:=\bigcup_{n=1}^{\infty}A_n.
\]
Since the sets are pairwise disjoint and $f\ge0$,
\[
f\mathbf1_{U_r}=\sum_{n=1}^{r}f\mathbf1_{A_n}
\uparrow f\mathbf1_U\qquad\text{pointwise on }S.
\]
The equality follows from $\mathbf1_{U_r}=\sum_{n=1}^{r}\mathbf1_{A_n}$; the monotone convergence follows from $U_r\uparrow U$.

Therefore,
\[
\begin{aligned}
\lambda\left(\bigcup_{n=1}^{\infty}A_n\right)&=\int_S f\mathbf1_U\,d\mu =\lim_{r\to\infty}\int_S f\mathbf1_{U_r}\,d\mu =\lim_{r\to\infty}\sum_{n=1}^{r}\int_S f\mathbf1_{A_n}\,d\mu =\sum_{n=1}^{\infty}\lambda(A_n).
\end{aligned}
\]
The second equality follows from MON (Theorem~5.2), the third from nonnegative linearity (Lemma~5.6), and the last from the definition of an infinite nonnegative sum. Hence $\lambda$ is countably additive and therefore is a measure on $(S,\Sigma)$.

Moreover, let $N\in\Sigma$ satisfy $\mu(N)=0$. Define
\[
h_r:=(f\wedge r)\mathbf1_N,\qquad r\in\mathbb N.
\]
Then $h_r\in(m\Sigma)^+$, $h_r\uparrow f\mathbf1_N$, and $0\le h_r\le r\mathbf1_N$. By monotonicity (HW~45),
\[
0\le\mu(h_r)\le\mu(r\mathbf1_N)=r\mu(N)=0.
\]
The equality uses the integral of a simple indicator. Applying MON once more gives
\[
\lambda(N)=\int_S f\mathbf1_N\,d\mu
=\lim_{r\to\infty}\mu(h_r)=0.
\]
Thus $\lambda\ll\mu$.

Therefore, the statement is proved. 

\subsection*{Theorem 5.8 --- Radon--Nikodym Theorem}%
\manualtoc{subsection}{\texorpdfstring{Theorem 5.8 --- Radon--Nikodym Theorem}{Theorem 5.8 --- Radon--Nikodym Theorem}}{studiabile-nav-3004}%

\studyheading{Statement.} Let $\mu$ and $\lambda$ be $\sigma$-finite measures on $(S,\Sigma)$ such that $\lambda\ll\mu$. Then there exists $f\in(m\Sigma)^+$ satisfying
\[
\lambda(A)=\int_S f\mathbf1_A\,d\mu,
\qquad\forall A\in\Sigma.
\]
The function
\[
f=: \frac{d\lambda}{d\mu}
\]
is called the Radon--Nikodym derivative of $\lambda$ with respect to $\mu$.

The density is unique up to $\mu$-almost-everywhere equality.

\par\noindent\studyheading{Recall.}%
\inlinetoc{subsubsection}{\texorpdfstring{Recall.}{Recall.}}{studiabile-nav-3005}%
\par\nobreak
 A measure $\nu$ is $\sigma$-finite if there exist measurable sets $(E_n)_{n\in\mathbb N}$ such that
\[
S=\bigcup_{n=1}^{\infty}E_n,
\qquad\nu(E_n)<\infty\quad\forall n\in\mathbb N.
\]
HW~53 constructs a measure $\lambda$ from a given nonnegative density $f$; the Radon--Nikodym theorem gives the converse under $\sigma$-finiteness and absolute continuity. As specified in the lecture notes, the theorem is stated without proof.
